\begin{textsample}{NEG Dim 2 – 2016\_c5 – Score -3.00 – 2016\_c5/enade\_2016\_fisio\_f3597e5bbb.txt}  \label{ex:f2_neg_035}
CAPÍTULO 5 – DISTRIBUIÇÃO DOS CONCEITOS 5.1 - PANORAMA NACIONAL DA DISTRIBUIÇÃO DOS CONCEITOS O Conceito Enade é um indicador de qualidade que avalia o desempenho dos estudantes a partir dos resultados do Enade.
Os valores possíveis vão de 1 ( pior situação ) a 5 ( melhor situação ).
As Unidades de Observação com apenas um ou sem nenhum concluinte participante não obtêm o Conceito Enade, ficando Sem Conceito ( SC ).
A Tabela 5.1 apresenta a quantidade e distribuição de cursos de Fisioterapia participantes do Enade/2016, por faixa de conceito e Grande Região. \textbf{Enfatiza} -se, mais uma vez, que a diferença entre os cursos tabulados neste capítulo e no capítulo 2 corresponde aos cursos sem conceito, ou seja àqueles, em princípio, sem alunos concluintes que houvessem participado da prova ou àqueles que tivessem tão somente um aluno concluinte19.
Observando -se os dados da Tabela 5.1, nota -se que, dos 463 cursos participantes, 212 ( 45,8 \% ) classificaram -se com conceito 3, o valor modal.
Esse foi também o conceito modal em todas as regiões : Norte ( 59,3 \% ), Nordeste ( 48,5 \% ), Sudeste ( 41,8 \% ), Sul ( 45,5 \% ) e Centro-Oeste ( 51,2 \% ).
O conceito 4 foi o segundo mais frequente em nível nacional ( 28,3 \%, correspondendo a 131 cursos ) e o conceito 2, o terceiro ( 16,8 \%, correspondendo a 78 cursos ).
Houve, ainda, 31 cursos ( 6,7 \% ) que receberam conceito 5 e seis cursos ( 1,3 \% ) que receberam conceito 1.
Dos 463 cursos de Fisioterapia, cinco ( 1,1 \% ) ficaram sem conceito ( SC ). 19 Estes cursos com somente um concluinte não foram considerados no capítulo 2.
A região Norte participou com 27 cursos ou 5,8 \% do total nacional.
Desses, 16 cursos receberam o conceito 3, o que equivale a 59,3 \% do total regional, além de ser o conceito modal na região, como já comentado.
A nove cursos ( 33,3 \% ) atribuiu -se conceito 2 e a dois cursos ( 7,4 \% ), conceito 4.
Nenhum curso ficou sem conceito ( SC ) ou com conceito 1 ou 5 nessa região.
A região Nordeste participou com 103 cursos ou 22,2 \% do total nacional.
Nessa região, 50 cursos ( 48,5 \% em termos regionais ) obtiveram conceito 3, o conceito modal para a região, como já comentado.
Os conceitos 2 e 4 foram atribuídos a 24 cursos cada, correspondendo a 23,3 \%, cada.
Ainda nessa região, cinco cursos ( 4,9 \% ) foram avaliados com o conceito 5.
Nenhum curso ficou sem conceito ( SC ) ou com conceito 1.
Dos 213 cursos participantes da região Sudeste ( 46,0 \% do total nacional ), o conceito 3 foi o conceito modal, como já anteriormente assinalado, tendo sido obtido por 89 cursos ( 41,8 \% ).
O conceito 1 foi atribuído a quatro cursos ( 1,9 \% ) e o conceito 2, a 27 cursos ( 12,7 \% ).
Receberam os conceitos 4 e 5, respectivamente, 76 cursos ( 35,7 \% ) e 14 cursos ( 6,6 \% ).
Nessa região, três cursos ( 1,4 \% ) ficaram sem conceito ( SC ).
Os 77 cursos da região Sul corresponderam a 16,6 \% do total nacional.
Como já indicado, houve predominância do conceito 3, atribuído a 45,5 \% dos cursos da região ( conceito modal ), ou seja, atribuído a 35 dos 77 cursos participantes na região Sul.
O conceito 2 foi atribuído a oito cursos ( 10,4 \% ), o conceito 4, a 22 cursos ( 28,6 \% ) e o conceito 5, a dez cursos ( 13,0 \% ).
Dois cursos ficaram sem conceito ( SC ) e nenhum curso recebeu conceito 1.
Já dos 43 cursos participantes na região Centro-Oeste ( 9,3 \% do total nacional ), 22 receberam conceito 3, o conceito modal ( 51,2 \% ), como já fora destacado.
A dois cursos foi atribuído conceito 1, e a outros dois, conceito 5 ( 4,7 \% cada ).
Os conceitos restantes, 2 e 4, foram atribuídos, respectivamente, a dez ( 23,3 \% ) e sete ( 16,3 \% ) cursos.
Nessa região, nenhum curso ficou sem conceito ( SC ).
O Gráfico 5.1 apresenta a distribuição do Conceito Enade segundo Grande Região : poligonais mais à esquerda representam Grandes Regiões com pior distribuição, e poligonais mais à direita, Grandes Regiões com melhores conceitos.
A região Norte ( linha azul ) apresenta os cursos com conceitos mais baixos, ao passo que as regiões Sudeste ( linha vermelha ) e Sul ( linha roxa ) apresentam as distribuições com valores maiores.
Cursos em IES situadas nas regiões Nordeste ( linha verde ) e Centro-Oeste ( linha preta ) apresentam uma situação intermediária. 5.2 - CONCEITOS POR CATEGORIA ADMINISTRATIVA E GRANDE REGIÃO A Tabela 5.2 apresenta a distribuição dos cursos participantes do Enade/2016 de Fisioterapia, por Categoria Administrativa, de acordo com os conceitos por eles alcançados, segundo as Grandes Regiões brasileiras.
Dos 463 cursos participantes, 58 ( 12,5 \% ) eram ministrados em Instituições Públicas e 405 ( 87,5 \% ), em Instituições Privadas.
De acordo com as informações da Tabela 5.2, em termos nacionais, entre Instituições Públicas e Privadas, observa -se que, dos 31 cursos avaliados com conceito 5, 18 eram oferecidos em IES Públicas e 13, em IES Privadas.
Dos 58 cursos participantes de IES Públicas, o conceito 4 foi o valor modal, atribuído a 24 cursos ( 41,4 \% ).
Entre os demais cursos de Instituições Públicas participantes, um curso ( 1,7 \% ) foi avaliado com conceito 1, cinco cursos ( 8,6 \% da categoria ) obtiveram conceito 2, dez cursos ( 17,2 \% ) receberam conceito 3 e, como já comentado, 18 cursos ( 31,0 \% ) receberam conceito 5.
Nesta categoria, nenhum curso ficou sem conceito ( SC ).
Na rede privada, o conceito modal foi 3, atribuído a 202 cursos, o correspondente a 49,9 \% dos 405 cursos da categoria.
Entre os demais cursos participantes, cinco ( 1,2 \% ) receberam conceito 1, e 73 ( 18,0 \% ), conceito 2.
O conceito 4 foi atribuído a 107 cursos ( 26,4 \% ) e o conceito 5, a 13 cursos ( 3,2 \% ).
Nesta Categoria Administrativa, cinco cursos ( 1,2 \% ) ficaram sem conceito ( SC ).
Na análise por região, observa -se que, na região Norte, as Instituições Privadas participaram com 23 cursos ( 82,5 \% do total regional ), dos quais 15 ficaram com conceito 3 e oito com conceito 2.
Nenhum ficou sem conceito ( SC ) ou recebeu os demais conceitos.
As Instituições Públicas participaram com quatro cursos na região Norte ( 14,8 \% do total regional ), dos quais dois obtiveram conceito 4, o valor modal.
Os conceitos 2 e 3 foram atribuídos a dois cursos ( um curso, cada ).
Nenhum curso ficou sem conceito ( SC ) ou com conceito 1 ou 5 nessa combinação de categoria e região.
Na região Nordeste, a rede privada concentrou 88 dos 103 cursos participantes, o equivalente a 85,4 \% do total da região.
Dentre os cursos oferecidos por IES Privadas no Nordeste, 50 cursos foram avaliados com conceito 3, o conceito modal.
Os demais foram avaliados com conceito 2, 24 cursos e conceito 4, 14 cursos.
Nessa combinação de categoria e região, nenhum curso ficou sem conceito ( SC ) ou com conceito 1 ou 5.
As Instituições Públicas da região Nordeste participaram com 15 cursos ( 14,6 \% do total da região ).
Desses, dez obtiveram conceito 4, o valor modal, e os outros cinco obtiveram conceito 5.
Nessa região nenhum curso ficou sem conceito ( SC ) ou com os demais conceitos.
Na região Sudeste, a proporção de cursos da rede privada, 90,6 \%, foi mais \textbf{elevada} do que a registrada nas demais regiões brasileiras, correspondendo a 193 dos 213 cursos participantes.
Nessa categoria e região, o conceito modal foi 3, atribuído a 85 cursos.
Três cursos ficaram sem conceito ( SC ).
Os demais foram avaliados com conceito 1 ( três cursos ), conceito 2 ( 25 cursos ), conceito 4 ( 68 cursos ) e conceito 5 ( nove cursos ).
Entre os 20 cursos oferecidos por Instituições Públicas na região Sudeste, o conceito modal foi 4, atribuído a oito cursos.
Os demais cursos receberam os conceitos 1 ( um curso ), 2 ( dois cursos ), 3 ( quatro cursos ) e 5 ( cinco cursos ).
Na região Sudeste, nesta categoria, nenhum curso ficou sem conceito ( SC ).
As Instituições Privadas concentraram 64 dos 77 cursos participantes da região Sul, 83,1 \% do total regional, a menor das proporções desta categoria entre as regiões.
Desses, 33 obtiveram conceito 3, o conceito modal.
Nesta combinação de Categoria Administrativa e Grande Região, os conceitos 2 e 4 foram atribuídos a sete e a 19 cursos, respectivamente.
Três cursos receberam conceito 5, e nenhum recebeu conceito 1.
Dois cursos ficaram sem conceito ( SC ).
As Instituições Públicas na região Sul participaram com 13 cursos, 16,9 \% dos cursos da região, a maior das proporções desta categoria nas regiões.
O conceito modal foi 5, atribuído a sete cursos.
Aos demais cursos foram atribuídos os conceitos 2 ( um curso ), 3 ( dois cursos ) e 4 ( três cursos ).
Nenhum curso da região Sul oferecido por IES Pública ficou sem conceito ( SC ) ou recebeu o conceito 1.
Na região Centro-Oeste, 37 dos 43 cursos participantes eram de Instituições Privadas ( 86,0 \% em termos regionais ).
Desses, 19 receberam conceito 3, o conceito modal.
Entre os demais, dois receberam conceito 1, nove receberam conceito 2, seis receberam conceito 4 e um recebeu conceito 5.
Nenhum curso ficou sem conceito ( SC ).
Quanto aos seis cursos oferecidos por Instituições Públicas na região Centro-Oeste ( 14,0 \% do total regional ), nota -se que três cursos foram alocados ao conceito 3, conceito modal.
Os outros três cursos ficaram com os conceitos 2, 4 e 5 ( um curso, cada ).
Nenhum curso ficou sem conceito ( SC ) ou com conceito 1.
O Gráfico 5.2 apresenta a distribuição do Conceito Enade segundo a Categoria Administrativa da IES.
Os cursos em IES Públicas ( linha azul ) apresentam uma poligonal mais à direita do que os em IES Privadas ( linha verde ) e, consequentemente, uma distribuição de Conceitos Enade com valores maiores. 5.3 - CONCEITOS POR ORGANIZAÇÃO ACADÊMICA E GRANDE REGIÃO Na Tabela 5.3, encontra -se a distribuição dos conceitos atribuídos aos cursos participantes do Enade/2016 na Área de Fisioterapia, por Organização Acadêmica, segundo as Grandes Regiões brasileiras.
Dos 463 cursos de Fisioterapia participantes, 189 eram oferecidos em Universidades, 92 em Centros Universitários e 182 em Faculdades.
Esta distribuição corresponde a, respectivamente, 40,8 \%, 19,9 \% e 39,3 \% dos cursos.
De acordo com os dados apresentados, de todos os 31 cursos avaliados com conceito 5, 24 eram vinculados a Universidades.
Os cursos vinculados a esse tipo de Organização Acadêmica tiveram o conceito 4 como conceito modal, atribuído a 75 cursos ( 39,7 \% ).
Os demais cursos vinculados a Universidades receberam os conceitos 1 ( um curso ), 2 ( 18 cursos ), 3 ( 69 cursos ) e conceito 5 ( 24 cursos, como já mencionado ).
Ficaram sem conceito ( SC ), dois cursos.
Entre os cursos em Centros Universitários, o conceito modal foi 3, atribuído a 49 cursos ( 53,3 \% ).
Dois cursos vinculados a esse tipo de Organização Acadêmica ficaram sem conceito ( SC ), e os outros cursos receberam os conceitos 1 ( um curso ), 2 ( oito cursos ), 4 ( 28 cursos ) e 5 ( quatro cursos ).
Dos 182 cursos mantidos por Faculdades, 94 ( 51,6 \% ) receberam conceito 3, o conceito modal.
O segundo conceito mais frequentemente atribuído aos cursos oferecidos por essa categoria de IES foi o conceito 2, recebido por 52 cursos.
Na sequência, 28 cursos obtiveram conceito 4 e quatro cursos, conceito 1.
O conceito 5 foi atribuído a três cursos, e um curso ficou sem conceito ( SC ).
Considerando -se separadamente as regiões brasileiras, verifica -se que, na região Norte, as Universidades concentraram seis dos 27 cursos participantes.
Estes cursos ficaram \textbf{distribuídos} nos conceitos 2, 3 e 4 ( dois cursos, cada ).
Os Centros Universitários da região Norte foram representados por seis cursos, aos quais todos foram atribuídos o conceito 3.
As Faculdades participaram com 15 cursos na região Norte que receberam os conceitos 2 ( sete cursos ) e 3 ( oito cursos, conceito modal ).
Na região Nordeste, as Universidades participaram com 24 dos 103 cursos da Área de Fisioterapia oferecidos na região, dos quais nenhum ficou sem conceito ( SC ) ou com conceito 1.
O conceito modal foi 4, atribuído a 13 cursos e os conceitos 3 e 5, a cinco cursos cada.
Um único curso obteve conceito 2.
Os Centros Universitários contaram com 12 cursos participantes na região Nordeste.
Esses obtiveram conceitos 3 ( nove cursos, conceito modal ) e 4 ( três cursos ).
Nenhum curso oferecido por Centros Universitários na região Nordeste ficou sem conceito ( SC ) ou recebeu os demais conceitos.
As Faculdades foram representadas por 67 cursos na região Nordeste, 36 dos quais obtiveram o conceito modal 3.
O conceito 2 foi atribuído a 23 cursos e o conceito 4, a oito cursos.
Na região Nordeste, nenhum curso oferecido por Faculdades ficou sem conceito ( SC ) ou recebeu conceito 1 ou 5.
Na região Sudeste, as Universidades concentraram 103 dos 213 cursos de Fisioterapia da região.
Entre os cursos em Universidades, nessa região, o conceito modal foi 4, atribuído a 41 cursos.
Os demais cursos receberam os conceitos 1 ( um curso ), 2 ( 12 cursos ), 3 ( 39 cursos ) e 5 ( dez cursos ).
Nenhum curso ficou sem conceito ( SC ).
Os Centros Universitários participaram com 52 cursos na região Sudeste, dos quais 24 obtiveram o conceito modal 3.
O conceito 1 foi atribuído a um único curso, o conceito 2, a seis cursos, o conceito 4, a 18 cursos e, por fim, o conceito 5, a um curso.
Dois cursos ficaram sem conceito ( SC ) nessa combinação de Organização Acadêmica e região.
Já as Faculdades foram representadas por 58 cursos na região Sudeste, dos quais um ficou sem conceito ( SC ).
Aos demais cursos foram atribuídos os conceitos 1 ( dois cursos ), 2 ( nove cursos ), 3 ( 26 cursos, o conceito modal ), 4 ( 17 cursos ) e 5 ( três cursos ).
Dos 77 cursos da região Sul, 43 eram de Universidades, para os quais o conceito modal foi 4, atribuído a 16 cursos.
Nesse tipo de organização, dois cursos ficaram sem conceito ( SC ) e nenhum recebeu conceito 1.
Os demais cursos receberam os conceitos 2 ( três cursos ), 3 ( 15 cursos ) e 5 ( sete cursos ).
Dos 14 cursos participantes de Centros Universitários da região Sul, nenhum ficou sem conceito ( SC ) ou recebeu conceito 1.
Os conceitos modais 3 e 4, foram atribuídos a cinco cursos, cada.
Os conceitos restantes, 2 e 5, foram atribuídos, respectivamente, a um e a três cursos.
Quanto aos 20 cursos vinculados a Faculdades na região Sul, 15 receberam o conceito modal 3.
Nenhum curso ficou sem conceito ( SC ).
O conceito 2 foi atribuído a quatro cursos, e o conceito 4, a um curso.
Nenhum curso recebeu conceito 1, tampouco obteve conceito 5.
Na região Centro-Oeste, 13 dos 43 cursos eram de Universidades.
Nesse tipo de organização, nenhum curso ficou sem conceito ( SC ) ou recebeu conceito 1 ou 2.
O conceito 3, conceito modal, foi atribuído a oito cursos, o conceito 4, a três cursos e o conceito 5, a dois cursos.
Os Centros Universitários da região Centro-Oeste contaram com oito cursos, que obtiveram os conceitos 2 ( um curso ), 3 ( cinco cursos, conceito modal ) e 4 ( dois cursos ).
Nenhum curso ficou sem conceito ( SC ) ou recebeu os conceitos restantes na categoria.
Dos 22 cursos oferecidos em Faculdades na região Centro-Oeste, 18 receberam os conceitos modais 2 e 3 ( nove cursos, cada ).
Aos demais cursos foram atribuídos os conceitos 1 e 4 ( dois cursos, cada ).
Nenhuma Faculdade da região Centro-Oeste ficou sem conceito ( SC ) ou com conceito 5.
O Gráfico 5.3 apresenta a distribuição do Conceito Enade segundo a Organização Acadêmica da IES.
Os cursos em Universidades ( linha azul ) apresentam uma poligonal mais à direita do que as demais, denotando uma melhor distribuição de Conceitos Enade.
No outro extremo, temos os cursos oferecidos em Faculdades ( linha vermelha ) com a poligonal mais à esquerda e, consequentemente, uma distribuição de Conceitos Enade com menores valores.
IES de Centros Universitários ( linha verde ) apresentam uma situação intermediária.

% matched lemmas: distribuir, elevado, enfatizar
\end{textsample}
